% GENERATED FILE. Edit canonical lessons, then run npm run generate:books.
% canonical-source-hash: fnv1a64:7584b976910a74d3
% canonical-lessons: JA-C339-bunbetsu-suru, JA-C339-hozon-suru, JA-C339-hogo-suru, JA-C339-toriireru, JA-C339-barabara

\chapter{Sorting and Keeping}
\label{ch:ja-sorting-and-keeping}
\hlchaptermodality{339}

\begin{chapteropening}
\textbf{By the end of this chapter:} \emph{I can explain how rubbish is sorted and things are kept.}

Complete the last lesson of chapter 339: I can explain how rubbish is sorted and things are kept.
\end{chapteropening}

\section[\texorpdfstring{\ja{ぶんべつする} (bunbetsu suru)}{ぶんべつする (bunbetsu suru)}]{\ja{ぶんべつする} (bunbetsu suru) — to sort rubbish; to separate waste}

\label{lesson:JA-C339-bunbetsu-suru}



\begin{quote}
Before the new one: say the Japanese for to burn something, then the Japanese for resources.
\end{quote}

\subsection*{You'll want to know: \ja{ぶんべつする}}
\textbf{\ja{ぶんべつする}} — \emph{bunbetsu suru} — \textquotedblleft{}to sort rubbish; to separate waste\textquotedblright{}.

\begin{quote}
\textbf{\ja{きそくが} \ja{あるので、ごみを} \ja{ぶんべつします。}} — \emph{kisoku ga aru node, gomi o bunbetsu shimasu} — \textquotedblleft{}There is a rule, so we sort our rubbish.\textquotedblright{}
\end{quote}

Polite \textbf{\ja{ぶんべつします}}, \emph{bunbetsushimasu}; a \ja{する} verb. Past (polite): \textbf{\ja{ははが} \ja{ごみを} \ja{ぶんべつしました。}}, \emph{haha ga gomi o bunbetsu shimashita}.

\begin{cousinweb}
\textbf{\ja{ぶんべつする}} (\emph{bunbetsu suru}) is built from a Sino-Japanese noun and the native verb \emph{suru}, \textquotedblleft{}to do\textquotedblright{}. The noun, \emph{bunbetsu}, is two Chinese characters read with their on'yomi, readings borrowed from Chinese: \emph{bun}, \textquotedblleft{}part\textquotedblright{} (Middle Chinese \emph{pjun}) and \emph{betsu}, \textquotedblleft{}separate\textquotedblright{} (\emph{bjet}) (KANJIDIC2; JMdict; Middle Chinese in Baxter's transcription).
\end{cousinweb}

\begin{grammarlens}[title={What you've built}]
One more word for giving reasons and opinions.
\end{grammarlens}

\subsection*{Your turn}
\begin{itemize}
\raggedright
  \item \emph{Say it:} \emph{bunbetsu suru}
  \item \emph{Say it:} \emph{bunbetsu suru}, once more
  \item \emph{Say it:} \emph{kisoku ga aru node, gomi o bunbetsu shimasu} — There is a rule, so we sort our rubbish.
  \item \emph{Recall:} say the Japanese for to burn something, then the Japanese for resources, then say \emph{bunbetsu suru} again
\end{itemize}

\begin{tcolorbox}[breakable,colback=teal!4,colframe=teal!35!black,title={Before you move on}]
What does \ja{ぶんべつする} mean? (\textquotedblleft{}To sort rubbish; to separate waste\textquotedblright{}.) Say it once more.
\end{tcolorbox}
\section[\texorpdfstring{\ja{ほぞんする} (hozon suru)}{ほぞんする (hozon suru)}]{\ja{ほぞんする} (hozon suru) — to keep; to preserve}

\label{lesson:JA-C339-hozon-suru}



\begin{quote}
Before the new one: say the Japanese for resources, then the Japanese for to sort rubbish.
\end{quote}

\subsection*{You'll want to know: \ja{ほぞんする}}
\textbf{\ja{ほぞんする}} — \emph{hozon suru} — \textquotedblleft{}to keep; to preserve\textquotedblright{}.

\begin{quote}
\textbf{\ja{あついので、たべものは} \ja{ひやして} \ja{ほぞんして} \ja{ください。}} — \emph{atsui node, tabemono wa hiyashite hozon shite kudasai} — \textquotedblleft{}It is hot, so keep food cool when you store it.\textquotedblright{}
\end{quote}

Polite \textbf{\ja{ほぞんします}}, \emph{hozonshimasu}; a \ja{する} verb. Past (polite): \textbf{\ja{みせが} \ja{たべものを} \ja{ほぞんしました。}}, \emph{mise ga tabemono o hozon shimashita}.

\begin{cousinweb}
\textbf{\ja{ほぞんする}} (\emph{hozon suru}) is built from a Sino-Japanese noun and the native verb \emph{suru}, \textquotedblleft{}to do\textquotedblright{}. The noun, \emph{hozon}, is two Chinese characters read with their on'yomi, readings borrowed from Chinese: \emph{ho}, \textquotedblleft{}preserve\textquotedblright{} (Middle Chinese \emph{pawX}) and \emph{zon}, \textquotedblleft{}exist\textquotedblright{} (\emph{dzwon}) (KANJIDIC2; JMdict; Middle Chinese in Baxter's transcription).
\end{cousinweb}

\begin{grammarlens}[title={What you've built}]
One more word for giving reasons and opinions.
\end{grammarlens}

\subsection*{Your turn}
\begin{itemize}
\raggedright
  \item \emph{Say it:} \emph{hozon suru}
  \item \emph{Say it:} \emph{hozon suru}, once more
  \item \emph{Say it:} \emph{atsui node, tabemono wa hiyashite hozon shite kudasai} — It is hot, so keep food cool when you store it.
  \item \emph{Recall:} say the Japanese for resources, then the Japanese for to sort rubbish, then say \emph{hozon suru} again
\end{itemize}

\begin{tcolorbox}[breakable,colback=teal!4,colframe=teal!35!black,title={Before you move on}]
What does \ja{ほぞんする} mean? (\textquotedblleft{}To keep; to preserve\textquotedblright{}.) Say it once more.
\end{tcolorbox}
\section[\texorpdfstring{\ja{ほごする} (hogo suru)}{ほごする (hogo suru)}]{\ja{ほごする} (hogo suru) — to safeguard; to protect (wildlife)}

\label{lesson:JA-C339-hogo-suru}



\begin{quote}
Before the new one: say the Japanese for to sort rubbish, then the Japanese for to preserve.
\end{quote}

\subsection*{You'll want to know: \ja{ほごする}}
\textbf{\ja{ほごする}} — \emph{hogo suru} — \textquotedblleft{}to safeguard; to protect (wildlife)\textquotedblright{}.

\begin{quote}
\textbf{\ja{どうぶつが} \ja{すくなく} \ja{なったので、ほごしています。}} — \emph{dōbutsu ga sukunaku natta node, hogo shite imasu} — \textquotedblleft{}The animals have become few, so they are being protected.\textquotedblright{}
\end{quote}

Polite \textbf{\ja{ほごします}}, \emph{hogoshimasu}; a \ja{する} verb. Past (polite): \textbf{\ja{まちが} \ja{もりを} \ja{ほごしました。}}, \emph{machi ga mori o hogo shimashita}.

\begin{cousinweb}
\textbf{\ja{ほごする}} (\emph{hogo suru}) is built from a Sino-Japanese noun and the native verb \emph{suru}, \textquotedblleft{}to do\textquotedblright{}. The noun, \emph{hogo}, is two Chinese characters read with their on'yomi, readings borrowed from Chinese: \emph{ho}, \textquotedblleft{}preserve\textquotedblright{} (Middle Chinese \emph{pawX}) and \emph{go}, \textquotedblleft{}protect\textquotedblright{} (\emph{huH}) (KANJIDIC2; JMdict; Middle Chinese in Baxter's transcription).
\end{cousinweb}

\begin{grammarlens}[title={What you've built}]
One more word for giving reasons and opinions.
\end{grammarlens}

\subsection*{Your turn}
\begin{itemize}
\raggedright
  \item \emph{Say it:} \emph{hogo suru}
  \item \emph{Say it:} \emph{hogo suru}, once more
  \item \emph{Say it:} \emph{dōbutsu ga sukunaku natta node, hogo shite imasu} — The animals have become few, so they are being protected.
  \item \emph{Recall:} say the Japanese for to sort rubbish, then the Japanese for to preserve, then say \emph{hogo suru} again
\end{itemize}

\begin{tcolorbox}[breakable,colback=teal!4,colframe=teal!35!black,title={Before you move on}]
What does \ja{ほごする} mean? (\textquotedblleft{}To safeguard; to protect (wildlife)\textquotedblright{}.) Say it once more.
\end{tcolorbox}
\section[\texorpdfstring{\ja{とりいれる} (toriireru)}{とりいれる (toriireru)}]{\ja{とりいれる} (toriireru) — to adopt an idea; to bring in}

\label{lesson:JA-C339-toriireru}



\begin{quote}
Before the new one: say the Japanese for to preserve, then the Japanese for to safeguard.
\end{quote}

\subsection*{You'll want to know: \ja{とりいれる}}
\textbf{\ja{とりいれる}} — \emph{toriireru} — \textquotedblleft{}to adopt an idea; to bring in\textquotedblright{}.

\begin{quote}
\textbf{\ja{しみんの} \ja{いけんを} \ja{とりいれた} \ja{ほうが} \ja{いいです。}} — \emph{shimin no iken o toriireta hō ga ii desu} — \textquotedblleft{}They should take the citizens' views on board.\textquotedblright{}
\end{quote}

Polite \textbf{\ja{とりいれます}}, \emph{toriiremasu}; an ichidan verb. Past (polite): \textbf{\ja{かいしゃが} \ja{とりいれました。}}, \emph{kaisha ga toriiremashita}.

\begin{cousinweb}
\textbf{\ja{とりいれる}} (\emph{toriireru}) is a native Japanese word (wago), inherited rather than borrowed from Chinese: it is built from \emph{tori}, \textquotedblleft{}taking (toru)\textquotedblright{} and \emph{ireru}, \textquotedblleft{}to put in\textquotedblright{}, and its kanji were attached later as kun'yomi (JMdict).
\end{cousinweb}

\begin{grammarlens}[title={What you've built}]
One more word for giving reasons and opinions.
\end{grammarlens}

\subsection*{Your turn}
\begin{itemize}
\raggedright
  \item \emph{Say it:} \emph{toriireru}
  \item \emph{Say it:} \emph{toriireru}, once more
  \item \emph{Say it:} \emph{shimin no iken o toriireta hō ga ii desu} — They should take the citizens' views on board.
  \item \emph{Recall:} say the Japanese for to preserve, then the Japanese for to safeguard, then say \emph{toriireru} again
\end{itemize}

\begin{tcolorbox}[breakable,colback=teal!4,colframe=teal!35!black,title={Before you move on}]
What does \ja{とりいれる} mean? (\textquotedblleft{}To adopt an idea; to bring in\textquotedblright{}.) Say it once more.
\end{tcolorbox}
\section[\texorpdfstring{\ja{ばらばら} (barabara)}{ばらばら (barabara)}]{\ja{ばらばら} (barabara) — scattered; in pieces}

\label{lesson:JA-C339-barabara}



\begin{quote}
Before the new one: say the Japanese for to safeguard, then the Japanese for to adopt an idea.
\end{quote}

\subsection*{You'll want to know: \ja{ばらばら}}
\textbf{\ja{ばらばら}} — \emph{barabara} — \textquotedblleft{}scattered; in pieces\textquotedblright{}.

\begin{quote}
\textbf{\ja{ふるいので、ほんが} \ja{ばらばらです。}} — \emph{furui node, hon ga barabara desu} — \textquotedblleft{}It is old, so the book is falling apart.\textquotedblright{}
\end{quote}

\begin{cousinweb}
\textbf{\ja{ばらばら}} (\emph{barabara}) is a native mimetic word, one of the Japanese words that picture a manner or a feeling in sound; JMdict marks it as mimetic. It is inherited, not borrowed, and its older origin is not settled (JMdict).
\end{cousinweb}

\begin{grammarlens}[title={What you've built}]
One more word for giving reasons and opinions.
\end{grammarlens}

\subsection*{Your turn}
\begin{itemize}
\raggedright
  \item \emph{Say it:} \emph{barabara}
  \item \emph{Say it:} \emph{barabara}, once more
  \item \emph{Say it:} \emph{furui node, hon ga barabara desu} — It is old, so the book is falling apart.
  \item \emph{Recall:} say the Japanese for to safeguard, then the Japanese for to adopt an idea, then say \emph{barabara} again
\end{itemize}

\begin{tcolorbox}[breakable,colback=teal!4,colframe=teal!35!black,title={Before you move on}]
What does \ja{ばらばら} mean? (\textquotedblleft{}Scattered; in pieces\textquotedblright{}.) Say it once more.
\end{tcolorbox}
